\section*{Chapitre \ref{ch:b2:plant-life-cycles} --- Cycles de vie et reproduction des plantes terrestres}

\begin{solution}{exo:b2:plant-life-cycles:1}
Gamétophyte~: la génération \omterm{def:b2:unicellular-diversity:unicellular}{pluricellulaire} \omterm{def:b2:meiosis-heredity:meiosis}{haploïde}, issue d'une spore
par mitoses, qui produit des gamètes par mitose. \omterm{def:b2:plant-life-cycles:cycles}{Sporophyte}~: la
génération \omterm{def:b2:meiosis-heredity:meiosis}{diploïde}, issue d'un zygote par mitoses, qui produit des
spores par \omterm{def:b2:meiosis-heredity:meiosis}{méiose}. Spore~: une cellule \omterm{def:b2:meiosis-heredity:meiosis}{haploïde} produite par \omterm{def:b2:meiosis-heredity:meiosis}{méiose}, qui
se développe sans fusionner. Gamète~: une cellule \omterm{def:b2:meiosis-heredity:meiosis}{haploïde} produite par
mitose (chez les plantes), qui doit fusionner avec une autre.
\end{solution}

\begin{solution}{exo:b2:plant-life-cycles:2}
Les cellules des feuilles sont à $2n = 1260$~: spore 630, cellule de
\omterm{def:b2:plant-life-cycles:fern}{prothalle} 630, spermatozoïde 630, oosphère 630, zygote 1260.
\end{solution}

\begin{solution}{exo:b2:plant-life-cycles:3}
Coussinet vert~: gamétophyte. Tige et capsule~: \omterm{def:b2:plant-life-cycles:cycles}{sporophyte}. Fronde de
\omterm{def:b2:plant-life-cycles:fern}{fougère}~: \omterm{def:b2:plant-life-cycles:cycles}{sporophyte}. \omterm{def:b2:plant-life-cycles:fern}{Prothalle}~: gamétophyte. Grain de \omterm{def:b2:plant-life-cycles:heterospory}{pollen}~:
gamétophyte mâle. Réserve nutritive de la graine (chez un pin)~:
gamétophyte femelle. Embryon de la graine~: le \omterm{def:b2:plant-life-cycles:cycles}{sporophyte} suivant.
\end{solution}

\begin{solution}{exo:b2:plant-life-cycles:4}
Haplophasique (\emph{Chlamydomonas})~: \omterm{def:b2:meiosis-heredity:meiosis}{méiose} dans le zygote, aussitôt.
Diplophasique (animaux, \emph{Fucus})~: la \omterm{def:b2:meiosis-heredity:meiosis}{méiose} produit les gamètes.
Haplodiplophasique (\omterm{def:b2:plant-life-cycles:moss}{mousses}, \omterm{def:b2:plant-life-cycles:fern}{fougères}, plantes à graines, nombreuses
algues)~: la \omterm{def:b2:meiosis-heredity:meiosis}{méiose} produit des spores dans les \omterm{def:b2:plant-life-cycles:fern}{sporanges} du
\omterm{def:b2:plant-life-cycles:cycles}{sporophyte}.
\end{solution}

\begin{solution}{exo:b2:plant-life-cycles:5}
$0.03/10^{-4} = \qty{300}{s}$, cinq minutes. En \qty{20}{min},
\qty{12}{cm}. La fécondation ne se fait qu'entre plantes distantes de
quelques centimètres~; les \omterm{def:b2:plant-life-cycles:moss}{mousses} forment donc des coussinets denses
où les deux sexes (ou les deux types d'organes) sont à portée.
\end{solution}

\begin{solution}{exo:b2:plant-life-cycles:6}
$v_s = 2\times(1.5\times 10^{-5})^2\times 1099\times 9.81/(9\times
1.8\times 10^{-5}) = \qty{0.030}{m/s}$, soit \qty{3}{cm/s}~; distance
$uh/v_s
= 2\times 0.5/0.03 = \qty{33}{m}$. Graine~: $2\times 20/0.8 =
\qty{50}{m}$ --- une unité plus lourde, libérée de plus haut et munie
d'une aile, va à peu près aussi loin.
\end{solution}

\begin{solution}{exo:b2:plant-life-cycles:7}
$200\times 40\times 64 = \num{512000}$ spores par fronde~; $5.1$
\omterm{def:b2:plant-life-cycles:fern}{prothalles}~; $0.1$ jeune \omterm{def:b2:plant-life-cycles:cycles}{sporophyte} par fronde.
\end{solution}

\begin{solution}{exo:b2:plant-life-cycles:8}
Une graine est un \omterm{def:b2:plant-life-cycles:heterospory}{ovule} conservé sur la plante mère~: le gamétophyte
femelle doit y être retenu, enveloppé et nourri, tandis que le
gamétophyte mâle doit voyager jusqu'à lui. Cela exige deux types de
spores aux destins différents. L'unique type de spore d'une plante
isosporée doit être libéré pour se développer en un gamétophyte
bisexué autonome~; la retenir retiendrait aussi la fonction mâle et
imposerait l'autofécondation à chaque plante, et plus rien n'aurait à
voyager.
\end{solution}

\begin{solution}{exo:b2:plant-life-cycles:9}
Spore~: $\tfrac43\pi(1.5\times 10^{-5})^3\times 1100 =
\qty{1.6e-11}{kg}$, soit \qty{16}{ng}~; énergie $0.5\times 1.6\times
10^{-8}\times 2\times 10^{4} = \qty{1.6e-4}{J}$. Graine~: \qty{5}{mg},
$0.5\times 5\times 10^{-3}\times 2\times 10^{4} = \qty{50}{J}$ ---
$3\times 10^{5}$ fois plus. La spore peut produire une petite lame
photosynthétique et rien d'autre avant d'avoir de la lumière~; la graine
peut produire une racine et une tige à l'obscurité, attendre toute une
saison, et survivre à la consommation ou à la dessiccation.
\end{solution}

\begin{solution}{exo:b2:plant-life-cycles:10}
$10^{11}/(10^{6}\times 20) = \num{5000}$ grains par mètre cube~; flux
$5000\times 10^{-7}\times 2 = \qty{1e-3}{}$ grain par seconde~; environ
\qty{1000}{s}, un quart d'heure, par grain. Le cône doit rester ouvert
et réceptif pendant des heures, voire des jours, et précisément au
moment où les cônes mâles libèrent leur \omterm{def:b2:plant-life-cycles:heterospory}{pollen} --- la pollinisation est
calée à la semaine près.
\end{solution}

\begin{solution}{exo:b2:plant-life-cycles:11}
La diploïdie masque les \omterm{def:b2:genome-diversification:mutation}{mutations} récessives délétères~; un corps
\omterm{def:b2:meiosis-heredity:meiosis}{diploïde} pourvu de tissus conducteurs peut être grand et vivre
longtemps, et la hauteur favorise à la fois la captation de la lumière
et la \omterm{thm:b2:plant-life-cycles:dispersal}{dispersion des spores}~; les spores produites en hauteur sur un
\omterm{def:b2:plant-life-cycles:cycles}{sporophyte} se dispersent dans l'air, alors que les gamètes doivent
nager. La génération \omterm{def:b2:meiosis-heredity:meiosis}{haploïde} persiste parce que \omterm{def:b2:meiosis-heredity:meiosis}{méiose} et fécondation
exigent une phase \omterm{def:b2:meiosis-heredity:meiosis}{haploïde}, et que le grain de \omterm{def:b2:plant-life-cycles:heterospory}{pollen} et le sac
embryonnaire sont les corps minimaux qui amènent les gamètes l'un vers
l'autre et nourrissent l'embryon.
\end{solution}

\begin{solution}{exo:b2:plant-life-cycles:12}
En cultivant des spores et en les suivant au microscope, il pouvait
établir qu'une spore se développe en un petit corps portant des organes
sexuels, que l'embryon naît de l'oosphère fécondée à l'intérieur de
l'\omterm{def:b2:plant-life-cycles:moss}{archégone}, que la plante porteuse de spores se développe à partir de
lui, et que les mêmes organes existent sous forme réduite dans l'\omterm{def:b2:plant-life-cycles:heterospory}{ovule}
des conifères~: l'alternance comme succession de corps et d'organes. Il
ne pouvait pas savoir ce qui distingue les deux générations à l'échelle
de la cellule. Les dénombrements de chromosomes montrèrent que les deux
corps diffèrent d'un facteur deux par le nombre de chromosomes, que la
\omterm{def:b2:meiosis-heredity:meiosis}{méiose} se situe à la formation des spores et la fécondation au niveau de
l'oosphère, et expliquèrent ainsi pourquoi l'alternance est
obligatoire.
\end{solution}

\begin{solution}{pb:b2:plant-life-cycles:1}
\textbf{1.} $10\times 200\times 40\times 64 = 5.1\times 10^{6}$
spores.
\textbf{2.} $\tfrac43\pi(1.5\times 10^{-5})^3\times 1100 =
\qty{1.6e-11}{kg}$ par spore~; récolte $8\times 10^{-5}$ kg, soit
\qty{80}{mg}.
\textbf{3.} $v_s = 2\times 2.25\times 10^{-10}\times 1099\times
9.81/(1.62\times 10^{-4}) = \qty{0.030}{m/s}$.
\textbf{4.} $x = 1.5\times 0.5/0.030 = \qty{25}{m}$~; disque de
$\pi\times 25^2 = \qty{2000}{m^2}$~; environ \num{2600} spores par
mètre carré.
\textbf{5.} Des courants ascendants de quelques centimètres par seconde
dépassent $v_s$~: une spore prise dans la turbulence reste en l'air
pendant des heures et parcourt des centaines de kilomètres~; une seule
spore fonde une population si son \omterm{def:b2:plant-life-cycles:fern}{prothalle} peut s'autoféconder (il
porte les deux types d'organes), ce qui explique que les îles lointaines
aient des flores de \omterm{def:b2:plant-life-cycles:fern}{fougères} riches.
\textbf{6.} Spore~: $0.5\times 1.6\times 10^{-8}\times 2\times 10^{4} =
\qty{1.6e-4}{J}$~; récolte~: \qty{820}{J}~; une graine de pin~:
\qty{50}{J} --- la récolte entière de spores équivaut à seize graines.
\textbf{7.} $5.1\times 10^{6}/2\times 10^{4} = 256$ \omterm{def:b2:plant-life-cycles:fern}{prothalles}.
\textbf{8.} $1 - 0.8^{6} = 1 - 0.26 = 0.74$.
\textbf{9.} $256\times 0.74\times 0.5\times 0.02 = 1.9$ \omterm{def:b2:plant-life-cycles:fern}{fougère} adulte
par saison.
\textbf{10.} $30\times 1.9 \approx 57$~: bien plus que l'unique
remplacement d'une population stable~; les chiffres de survie sont donc
optimistes --- dans un habitat saturé, la plupart des jeunes
\omterm{def:b2:plant-life-cycles:cycles}{sporophytes} meurent de l'encombrement et de l'ombre.
\textbf{11.} $10^{-3}/10^{-4} = \qty{10}{s}$. De nombreuses \omterm{def:b2:plant-life-cycles:fern}{fougères}
font mûrir \omterm{def:b2:plant-life-cycles:moss}{anthéridies} et \omterm{def:b2:plant-life-cycles:moss}{archégones} à des moments différents, ou
libèrent une hormone (l'anthéridiogène) qui rend mâles les jeunes
\omterm{def:b2:plant-life-cycles:fern}{prothalles} voisins, de sorte que les spermatozoïdes proviennent d'un
autre individu.
\textbf{12.} Il est épais d'une seule cellule, sans racines, sans
protection, a besoin d'eau liquide pour la fécondation et vit quelques
semaines~: presque toute la mortalité du cycle porte sur lui (une spore
sur $2\times 10^{4}$ en devient un~; un quart d'entre eux ne connaissent
jamais de nuit humide).
\textbf{13.} $v_s = 2\times 6.25\times 10^{-10}\times 599\times
9.81/(1.62\times 10^{-4}) = \qty{0.045}{m/s}$.
\textbf{14.} $3\times 15/0.045 = \qty{1000}{m}$.
\textbf{15.} $10^{11}/(2\times 10^{7}) = \num{5000}$ par mètre cube.
\textbf{16.} $5000\times 10^{-7}\times 2 = \qty{1e-3}{s^{-1}}$~: 3.6 par
heure~; le premier arrive au bout d'environ \qty{1000}{s}, un quart
d'heure.
\textbf{17.} À cinq kilomètres sous le vent, le nuage s'est étalé
latéralement et vers le haut, et la plupart des grains se sont déposés
(une portée de \qty{1}{km} pour \qty{15}{m} de hauteur)~: la
concentration y est inférieure de plusieurs ordres de grandeur, et un
\omterm{def:b2:plant-life-cycles:heterospory}{ovule} peut attendre des jours avant de recevoir un grain~; le propre
\omterm{def:b2:plant-life-cycles:heterospory}{pollen} de l'arbre donne surtout des graines autofécondées, qui avortent.
La pollinisation par le vent exige un nuage dense, d'où les peuplements.
\textbf{18.} $10^{14}$ grains pour $2\times 10^{6}$ \omterm{def:b2:plant-life-cycles:heterospory}{ovules}~: $5\times
10^{7}$ grains par \omterm{def:b2:plant-life-cycles:heterospory}{ovule}.
\textbf{19.} $2\times 10^{6}$ \omterm{def:b2:plant-life-cycles:heterospory}{ovules} $\times 0.6\times 0.8 = 9.6\times
10^{5}$ graines~; \qty{4.8}{kg}~; \qty{48}{MJ}.
\textbf{20.} $3\times 20/0.8 = \qty{75}{m}$~: le \omterm{def:b2:plant-life-cycles:heterospory}{pollen} va à un
kilomètre, la graine à quelques hauteurs d'arbre~; les gènes voyagent
par le \omterm{def:b2:plant-life-cycles:heterospory}{pollen}, la plante par la graine.
\textbf{21.} $9.6\times 10^{5}\times 0.05\times 0.01 = 480$ descendants
adultes, contre 57 pour la \omterm{def:b2:plant-life-cycles:fern}{fougère}.
\textbf{22.} Pin~: $4.8\times 10^{7}/480 = \qty{100}{kJ}$ par
descendant adulte~; \omterm{def:b2:plant-life-cycles:fern}{fougère}~: $820/1.9 = \qty{430}{J}$ par descendant
adulte, deux cents fois moins.
\textbf{23.} Toutes deux contiennent moitié de matière sèche à
\qty{20}{kJ/g}~; avec un métabolisme de base de, disons,
\qty{0.5}{mW} par gramme sec, toutes deux tiennent $10^{4}\,
\text{J/g}/(5\times 10^{-4}\,\text{W/g}) = 2\times 10^{7}$ s, environ
huit mois~: ce n'est pas la durée de la dormance que paie l'énergie de
la graine. Elle paie ce qui suit la germination --- une racine et une
tige construites à l'obscurité à partir des réserves, des semaines
avant que la plantule ne se nourrisse elle-même --- et, avec le
tégument, la protection en attendant~; les \qty{1.6e-4}{J} de la spore
construisent quelques cellules qui doivent faire aussitôt la
photosynthèse.
\textbf{24.} Graine~: fécondation sans eau~; un embryon pourvu de
réserves, protégé et dormant, qui s'installe dans des milieux secs ou
saisonniers~; dispersion par les ailes, les fruits et les animaux. La
stratégie de la \omterm{def:b2:plant-life-cycles:fern}{fougère} l'emporte dans les habitats humides, ombragés
et stables, et pour coloniser des terrains lointains ou nouvellement
ouverts, où des propagules bon marché, légères et nombreuses comptent
davantage que les réserves.
\textbf{25.} \omterm{def:b2:plant-life-cycles:fern}{Fougère}~: $5.1\times 10^{6}$ spores par saison pour 1.9
adulte, soit $2.7\times 10^{6}$ spores par descendant adulte,
\qty{430}{J} chacun~; pin~: $9.6\times 10^{5}$ graines pour 480 adultes,
soit \num{2000} graines par descendant adulte, \qty{100}{kJ} chacun.
\end{solution}
